\documentclass{beamer}
\usetheme{Madrid}
\usepackage[utf8]{inputenc}
\usepackage[brazil]{babel}

\title{Relatório do Trabalho Prático 2}
\author{Felipe Kelemen \\ Gustavo Bastos \\ Octavio Carneiro}
\date{\today}

\begin{document}

\begin{frame}
  \titlepage
\end{frame}

\begin{frame}{Visão geral}
\begin{itemize}
  \item Arquitetura: \textbf{Agent}, \textbf{Manager}, \textbf{MIB}
  \item Objetivo: coletar e expor métricas do sistema de forma simples e eficiente
\end{itemize}
\end{frame}

\begin{frame}{Fluxo de comunicação}
\begin{enumerate}
  \item Cada \textbf{agent} coleta métricas localmente e expõe via TCP
  \item O \textbf{manager} conecta nos agents e solicita \texttt{GET <oid>} ou \texttt{GET .}
  \item Resposta: string no formato \texttt{oid=valor;oid=valor;...}
\end{enumerate}
\end{frame}

\begin{frame}{Componentes principais}
\begin{itemize}
  \item \textbf{MetricsCollector}: leitura de /proc e /sys, calcula deltas
  \item \textbf{SocketServer + WorkerPool}: aceita conexões e processa clientes
  \item \textbf{MIB (árvore)}: registra getters por OID e serializa respostas
  \item \textbf{Manager}: cliente síncrono que consolida métricas
\end{itemize}
\end{frame}

\begin{frame}{Métricas coletadas}
\begin{itemize}
  \item Uso de CPU (\%)
  \item Uso de memória (\%)
  \item Uptime (s)
  \item Taxas de rede (bytes/s)
\end{itemize}
\end{frame}

\begin{frame}{Principais decisões de projeto}
\begin{itemize}
  \item Protocolo textual simples facilita depuração
  \item Consulta \texttt{GET .} retorna todas as folhas — menor overhead de protocolo
  \item Manager síncrono: simples, porém escala linearmente com número de agents
\end{itemize}
\end{frame}

\begin{frame}{Resultados resumidos}
\begin{itemize}
  \item Testes com falhas de agentes: reconexão automática funciona
  \item Escalabilidade: até 1000 agents testados; tempo de consulta cresce linearmente
  \item Latência: resposta por agent muito baixa; total cresce com número de agentes
\end{itemize}
\end{frame}

\begin{frame}{Limitações e melhorias}
\begin{itemize}
  \item Tornar o \textbf{manager} assíncrono para consultas paralelas
  \item Adicionar segurança (autenticação/criptografia)
  \item Expandir MIB para operações adicionais (SET, GETNEXT, etc.)
\end{itemize}
\end{frame}


\end{document}
